\chapter{Propuesta de Validación del Sistema}

El presente capítulo define la estrategia propuesta para validar el funcionamiento de la
estación inteligente de control ESD. La validación se plantea a partir de pruebas funcionales
por subsistema, pruebas de integración y pruebas de operación en el laboratorio TEM
de Intel Costa Rica. El objetivo principal de esta etapa es comprobar que el prototipo
cumpla con los requerimientos funcionales definidos en el proyecto y que permita mejorar
la trazabilidad de las pruebas ESD mediante identificación de usuario, almacenamiento de
registros, monitoreo ambiental e interacción mediante interfaz gráfica.

\section{Estrategia general de validación}

La validación del prototipo se dividirá en tres niveles. El primer nivel corresponde a la
validación individual de los subsistemas electrónicos y de software. El segundo nivel
corresponde a la validación de integración, donde se verificará la comunicación entre los
módulos y la correcta ejecución del flujo completo de operación. Finalmente, el tercer nivel
corresponde a la validación operativa del sistema en el laboratorio TEM de Intel Costa
Rica, bajo las condiciones reales disponibles durante el periodo de ejecución del proyecto.

La validación se enfocará en comprobar las siguientes funciones principales:

\begin{itemize}
    \item Verificación de pulsera antiestática mediante el probador HZR-171.
    \item Identificación de usuario mediante lectura de código de barras.
    \item Identificación complementaria mediante reconocimiento facial.
    \item Visualización de instrucciones y resultados en la interfaz gráfica táctil.
    \item Medición de temperatura y humedad relativa.
    \item Almacenamiento local de registros de operación.
    \item Generación de reportes históricos.
\end{itemize}

\section{Validación del subsistema de verificación ESD}

El subsistema de verificación ESD se validará comprobando la integración del probador
HZR-171 con el controlador principal. Para esto, se verificarán las señales de salida del
probador correspondientes a los estados de prueba y su correcta interpretación por parte
del sistema.

Durante esta etapa se realizarán pruebas con condiciones representativas de operación,
verificando que el resultado obtenido por el probador sea mostrado correctamente en la
interfaz gráfica y almacenado en el registro del sistema. La prueba será considerada
satisfactoria si el sistema identifica correctamente los estados generados por el probador
y los asocia al usuario correspondiente.

\section{Validación del subsistema de identificación}

El subsistema de identificación se validará mediante dos mecanismos: lectura de código de
barras y reconocimiento facial. Para la lectura de código de barras, se utilizarán códigos
asociados a usuarios de prueba, verificando que el lector GM67 pueda capturar el
identificador y transmitirlo correctamente al controlador principal.

Para el reconocimiento facial, se registrarán usuarios de prueba y se evaluará la capacidad
del sistema para reconocerlos bajo condiciones controladas de iluminación y posición. Este
mecanismo se considerará complementario, por lo que la identificación mediante código de
barras será el método principal de autenticación del usuario.

\section{Validación del monitoreo ambiental}

El monitoreo ambiental se validará mediante la lectura de temperatura y humedad relativa
utilizando el sensor SHT31-D. La validación buscará comprobar que el sistema 
sea capaz de adquirir, mostrar y almacenar
las variables ambientales junto con cada registro de prueba ESD. Esta información permitirá
relacionar las condiciones del ambiente con los eventos registrados por la estación.

\section{Validación de la interfaz gráfica}

La interfaz gráfica táctil se validará mediante pruebas de navegación y operación del flujo
de usuario. Se comprobará que la pantalla muestre instrucciones claras, información del
usuario, estado del sistema, resultado de la prueba ESD y variables ambientales.

La interfaz será considerada funcional si permite ejecutar el flujo de prueba sin necesidad
de intervención externa y si presenta los resultados de manera clara para el usuario.

\section{Validación del almacenamiento y reportes}

El subsistema de almacenamiento se validará mediante la generación de registros locales
en memoria microSD. Cada registro deberá incluir, como mínimo, el identificador del
usuario, fecha, hora, resultado de la prueba ESD, temperatura y humedad relativa.

Además, se verificará la generación de reportes históricos en formato CSV. La validación
será satisfactoria si los registros se almacenan sin pérdida de información y si los reportes
pueden ser consultados posteriormente desde la memoria del sistema.

\section{Validación integral del sistema}

La validación integral consistirá en ejecutar el flujo completo de operación del prototipo.
El procedimiento general propuesto será el siguiente:

\begin{enumerate}
    \item El usuario se identifica mediante código de barras o reconocimiento facial.
    \item El sistema muestra la información del usuario en la interfaz gráfica.
    \item El usuario realiza la prueba de pulsera antiestática.
    \item El sistema lee el resultado del probador HZR-171.
    \item El sistema mide temperatura y humedad relativa.
    \item El resultado de la prueba se muestra en pantalla.
    \item El sistema almacena el registro localmente.
    \item El sistema genera o actualiza el reporte histórico correspondiente.
\end{enumerate}

Esta prueba permitirá comprobar la interoperabilidad de los subsistemas y verificar que el
prototipo cumpla con la funcionalidad principal planteada en los objetivos del proyecto.

\section{Criterios de aceptación}

La Tabla~\ref{tab:criterios_validacion} resume los criterios de aceptación propuestos para
la validación del sistema.

\begin{table}[H]
\centering
\caption{Criterios de aceptación para la validación del sistema.}
\label{tab:criterios_validacion}
\begin{tabularx}{\textwidth}{p{4.2cm} X}
\toprule
\textbf{Subsistema} & \textbf{Criterio de aceptación} \\
\midrule
Verificación ESD &
El sistema debe identificar correctamente los estados entregados por el probador HZR-171
y mostrarlos en la interfaz gráfica. \\

Identificación por código de barras &
El sistema debe leer correctamente los identificadores de usuario utilizados durante las
pruebas funcionales. \\

Reconocimiento facial &
El sistema debe reconocer usuarios previamente registrados bajo condiciones controladas
de iluminación y posición. \\

Monitoreo ambiental &
El sistema debe adquirir, mostrar y almacenar temperatura y humedad relativa durante la
operación del prototipo. \\

Interfaz gráfica &
La interfaz debe guiar al usuario durante el proceso de prueba y mostrar el resultado final
de forma clara. \\

Almacenamiento local &
Cada prueba ejecutada debe generar un registro local con identificador de usuario, fecha,
hora, resultado ESD y variables ambientales. \\

Generación de reportes &
El sistema debe generar reportes históricos en formato CSV a partir de los registros
almacenados. \\

Integración del sistema &
El flujo completo de operación debe ejecutarse correctamente desde la identificación del
usuario hasta el almacenamiento del resultado. \\
\bottomrule
\end{tabularx}
\end{table}

\section{Validación operativa en el laboratorio}

Finalmente, se realizarán pruebas de operación en el laboratorio TEM de Intel Costa
Rica, de acuerdo con las condiciones y permisos disponibles durante el desarrollo del
proyecto. Esta etapa permitirá observar el comportamiento del sistema en un entorno real
de laboratorio y detectar posibles ajustes necesarios en la interfaz, el montaje físico, el
flujo de usuario o el almacenamiento de datos.

La validación operativa no pretende certificar el prototipo como un sistema industrial de
control ESD, sino comprobar su funcionamiento como una plataforma académica y de
investigación orientada a mejorar la trazabilidad de pruebas ESD en entornos controlados.



